\chapter*{Abstract}

This study aims to examine the impact of participation in the Smart Farmer project on farmers’ skills. The study utilizes data from both participating and non-participating farmers during the period 2022–2023. The analysis employs econometric techniques combined with machine learning approaches, including Propensity Score Matching (PSM) and Difference-in-Differences (DID), to address selection bias and estimate causal effects.

The findings indicate that although participants in the program exhibit higher skill levels compared to non-participants both before and after the intervention, the DID results show that participation does not significantly improve farmers’ skills relative to the control group. Furthermore, machine learning and SHAP analysis reveal that key determinants of skill improvement include baseline skills, technological skills, education level, and infrastructure conditions such as irrigation access.

The results highlight substantial heterogeneity in outcomes across farmers and suggest that one-size-fits-all training programs may be ineffective. Therefore, targeted and personalized policy interventions are necessary to enhance the effectiveness of agricultural development programs.

\vspace{0.5cm}
\textbf{Keywords:} Smart Farmer, Farmer Skills, Impact Evaluation, Propensity Score Matching, Machine Learning